\section{논의}

\paragraph{해석.}
가장 방어 가능한 해석은 예측 지연 보정이다. 역방향 스텝이 적으면 연속된 잡음 제거 분포 또는
전이 분포가 크게 이동할 수 있다. 최근 Fisher--Rao 호를 연장하면 다음 예측을 앞서 추정하여 품질을
개선할 수 있다. 격자가 촘촘해질수록 원시 예측의 변화가 작아지므로 외삽의 이점은 줄어들고,
확률적 오차나 모델 오차를 증폭할 수 있다.

\paragraph{MDLM과 SEDD의 차이.}
MDLM은 흡수 전이 이전의 정답 토큰 예측에 모멘텀을 적용하고 완료된 토큰을 고정한다. SEDD는
확률적 상태에 의존하는 한 스텝 전이에 모멘텀을 적용한다. 후자의 대상은 시간에 따라 덜 매끄러울 수
있으며, 이는 곡선 외삽이라는 해석을 약화시키고 보존된 결과의 불안정한 경향과도 일치한다.

\paragraph{연속 기하 모델과의 관계.}
본 방법은 Fisher-Flow 및 RDLM과 동일한 Fisher--Rao 매장을 사용하지만, 그 동역학을 사용하지는
않는다. 진정한 리만 다단계 솔버에는 연속 상태 $X_t$, 잘 정의된 접벡터장, $T_{X_{i-1}}\mathcal{M}$과
$T_{X_i}\mathcal{M}$ 사이의 평행이동, 그리고 지수사상에 의한 상태 갱신이 필요하다. 사전학습된
이산 MDLM 또는 SEDD를 위해 그러한 솔버를 확립하려면 범주형 예측과 연속 벡터장을 연결하는 새로운
유도가 필요하다.
